\section{Test categories and the De Morgan arch}\label{sec:test}

This section takes up issue~\#187 of the edition: the questions that folders 108, 118 and 134-2 leave open on test categories, read as bridges. We fix a fundamental localizer $W$ (Grothendieck's \emph{localisateur fondamental}) on the category $\mathrm{Cat}$ of small categories; the usual weak equivalences, those whose nerve is a weak equivalence of simplicial sets, are the smallest one. For a small category $A$, $i_A\colon\widehat A\to\mathrm{Cat}$ sends a presheaf $X$ to its category of elements $A/X$, and $i_A^*(C)(a) = \Hom(A/a, C)$. We use the definitions of \cite[\S4.1]{cisinski-308}, which follow \emph{Pursuing Stacks}: $A$ is a \emph{weak test category} if $i_A^*(C)$ is $W$-aspherical for every small $C$ with a final object (Déf.~4.1.8); a \emph{local test category} if every $A/a$ is a weak test category (Déf.~4.1.12); a \emph{test category} if it is both. The Lawvere object $L_A = i_A^*(\Delta_1)$ is the presheaf of sieves, and Grothendieck's criterion reads: $A$ is a local test category iff $L_A$ is locally aspherical, that is, iff $A/(a\times L_A)$ is aspherical for every object $a$ (\cite[Th.~4.1.19]{cisinski-308}, \cite[Th.~6.6]{maltsiniotis-2001}). This is a bridge in the sense of \S\ref{sec:recall}, with a homotopical arch: a property of the topos $\widehat A$, that its sheaf of truth values is contractible, read on the category $A$.

\subsection{The questions and what is known}

The table lists the questions on test categories, weak, local and strict test categories, asphericity and separating intervals that the three folders leave open, by folder and page of the transcription. Folder 134-2 holds only the plans and tables of contents of \emph{Pursuing Stacks} (its text is printed); for it we give the section the table announces and read the section in the arXiv edition \cite{ps-arxiv}. \q{Ours} marks a step that is the edition's; \q{open} means not found in the sources we read, which are \cite{cisinski-308,maltsiniotis-2001,ps-arxiv,hubert-thesis,hubert-dk}. We did not consult Jardine's papers, Cisinski's \emph{Higher categories and homotopical algebra}, nor the printed Astérisque~301.

\begin{center}\scriptsize\emergencystretch=3em
\begin{tabularx}{\linewidth}{>{\raggedright\arraybackslash}p{1.25cm}>{\raggedright\arraybackslash}X>{\raggedright\arraybackslash}p{5.2cm}}
\toprule
folder, page & question & status\\
\midrule
108, p.~2, q.~2 & are $W$ and $W_A$ saturated? & settled: every fundamental localizer is strongly saturated \cite[Prop.~4.2.4]{cisinski-308}; so is $W_A$ (118, pp.~39--40)\\
108, p.~2, q.~3 & is $(\mathrm{Cat})$ a model category? & settled: Thomason's structure \cite[Th.~5.2.12, 5.2.15]{cisinski-308}, for every accessible $W$\\
108, p.~2, q.~4 & exactness of $\mathrm{Hot}_W$ and of $\mathrm{Cat}\to\mathrm{Hot}_W$, $\widehat A\to\mathrm{Hot}_W$ & partly: finite sums and products \cite[\S9]{maltsiniotis-2001}, filtered colimits \cite[Cor.~4.2.22]{cisinski-308}, products and $\widehat A\to\mathrm{Hot}$ iff $A$ totally aspherical \cite[Prop.~4.3.2]{cisinski-308} (with the correction of row \nolinkurl{118-groupoid-products-not-totally-aspherical})\\
108, p.~2, q.~5 & when is $\widehat A$ a closed model category? & settled: a model structure with cofibrations the monomorphisms and weak equivalences $W_A$ exists iff $A$ is a local test category \cite[introduction, Th.~1.4.3, Cor.~4.2.18]{cisinski-308}\\
108, p.~2, q.~6 & is $\Delta^{\circ}$ a test category? (\q{no}) & settled, ours: $\Delta^\circ$ has an initial object, so it is neither a weak nor a local test category (Theorem~\ref{th:ore} below); the test case is stated without argument in \cite[\S88]{ps-arxiv}\\
108, p.~2, q.~6a & the variants $\Delta_f$, cubes & $\Delta_f$ is a weak test category, not a test category \cite[\S8]{maltsiniotis-2001}; the cube category is a test category \cite[Cor.~8.4.13]{cisinski-308}; the other cubical signs are not identified on the page\\
108, p.~2, q.~7, 17 & invariants of $\widehat A$, Dold--Puppe & partly: homotopical Dold--Kan correspondences for the test categories satisfying a Whitehead condition, $\Theta$ included \cite{hubert-thesis,hubert-dk}\\
108, p.~3, q.~13 & aspherical with 0-connected topos $\Rightarrow$ totally aspherical? & partly: total asphericity is the asphericity of the $a\times b$ \cite[4.3.1--4.3.2]{cisinski-308}, \cite[Prop.~7.1]{maltsiniotis-2001}; a totally aspherical weak test category is a strict test category \cite[Prop.~1.17]{hubert-dk}; the 0-connected form of the page: open\\
108, p.~4, q.~16 & for a contractor $A$, are $i_A$, $j_A$ inverse up to homotopisms; contractibility through $A/X$; $i_A(X\times Y)\to i_AX\times i_AY$? & up to weak equivalences: yes, since a nontrivial contractor is a strict test category (\cite[\S88]{ps-arxiv}, \cite[\S7]{maltsiniotis-2001}) and \cite[Th.~4.1.26, Prop.~4.3.2]{cisinski-308}; up to homotopisms: open\\
118, p.~14, pp.~18--22 & pushouts along immersions in $(\mathrm{Cat})$ & the page answers \q{non} on p.~22; Cisinski's pushout criteria \cite[\S5.1]{cisinski-308} (row \nolinkurl{118-pushout-along-cofibrations-in-cat})\\
118, pp.~23--25 & calculus of fractions for $\mathrm{Hot}_A$ with a separating interval; which class $\mathfrak M$ in $(\mathrm{Cat})$? & in $\widehat A$ the page needs $A$ local test, as \cite[Th.~4.1.19]{cisinski-308}; the version in $(\mathrm{Cat})$: open\\
118, p.~14 (1) & $W = W^{\circ}$? & settled \cite[Sorites~3.3.8 (7)]{cisinski-308}\\
118, pp.~26--30 & arbitrary sums in $\mathrm{Hot}(W)$ & settled for accessible $W$, ours from \cite[3.3.8 (1), Th.~5.2.15]{cisinski-308}\\
118, p.~36 & is an infinite product of aspherical objects aspherical? & no, ours (Proposition~\ref{prop:products})\\
134-2, p.~5, \S\S31--33 & a criterion for test categories, necessity of T1--T3 & settled \cite[Th.~4.1.19]{cisinski-308}, \cite[Th.~6.6]{maltsiniotis-2001}\\
134-2, p.~6 \S43, p.~9 \S98 & $\Delta^f$ weak test; \q{perplexities about its being aspheric} & settled \cite[\S8]{maltsiniotis-2001}\\
134-2, p.~9, \S84 & the inclusion $f_!(M_{\mathrm{as}})\subset M'_{\mathrm{as}}$ & left unanswered in \cite[\S88, 1)]{ps-arxiv}; open\\
134-2, p.~9, \S86 & does a notion of map of asphericity structures exist? & open\\
134-2, p.~9, \S90 & existence of (perfect) test functors into a modelizer & open in general; for $\Delta\to\mathrm{Cat}$, Thomason's $c\,\mathrm{Sd}^2$ gives a Quillen equivalence \cite[Th.~5.2.12]{cisinski-308}\\
134-2, p.~9, \S91 & conditions on $A$ for the homotopy types of finite type; test categories with boundary operations & open; for homology, \cite{hubert-dk}\\
\bottomrule
\end{tabularx}
\end{center}

\subsection{The De Morgan obstruction}

By the De Morgan arch (arch~B of \cite{bridges-generator}), $\widehat A$ is De Morgan iff $A$ satisfies the Ore condition. The Lawvere object sees the Ore condition directly.

\begin{lemma}\label{lem:ore-sieves}
$A$ satisfies the Ore condition iff, for every arrow $f\colon b\to a$ and every nonempty sieve $S$ on $a$, the sieve $f^*S$ is nonempty. Equivalently, the empty sieves form a direct summand of $L_A$: $L_A = L_A^{\emptyset}\sqcup L_A^{+}$, where $L_A^{+}(a)$ is the set of nonempty sieves on $a$.
\end{lemma}

\begin{proof}
If $S$ contains $g\colon c\to a$ and the cospan $(f,g)$ completes to $f h = g k$, then $h\in f^*S$. Conversely, given a cospan $(f,g)$, apply the condition to the sieve generated by $g$: an arrow $h$ of $f^*S$ is an $h$ with $fh = gk$ for some $k$. The empty sieves always form a sub-presheaf; their complement is one iff no nonempty sieve pulls back to the empty one.
\end{proof}

\begin{theorem}\label{th:ore}
Let $W$ be a fundamental localizer other than the two coarse ones, and $A$ a nonempty small category satisfying the Ore condition, that is, with $\widehat A$ De Morgan. Then $A$ is neither a weak $W$-test category nor a local $W$-test category. More generally, if some $A/X$ satisfies the Ore condition, $A$ is not a local $W$-test category.
\end{theorem}

\begin{proof}
Every non-coarse fundamental localizer is contained in the localizer $W_0$ of the functors that induce a bijection on $\pi_0$ \cite[Prop.~9.3.2]{cisinski-308}, so a $W$-aspherical category is connected. By Lemma~\ref{lem:ore-sieves}, $A/L_A$ is the disjoint union of $A/L_A^{\emptyset}\simeq A$ and $A/L_A^{+}$, which contains the maximal sieves; both are nonempty, so $L_A$ is not aspherical, while a weak test category needs it to be (Déf.~4.1.8 with $C=\Delta_1$). Likewise $a\times L_A = (a\times L_A^{\emptyset})\sqcup(a\times L_A^{+})$, both pieces containing a section over $a$, so $L_A$ is not locally aspherical and $A$ is not a local test category \cite[Th.~4.1.19]{cisinski-308}. The last assertion follows, since $A/X$ is a local test category as soon as $A$ is \cite[Prop.~4.1.14 (v)]{cisinski-308}, and a slice of a nonempty category is nonempty.
\end{proof}

\begin{corollary}\label{cor:ore}
For the usual weak equivalences, none of the following is a weak test category or a local test category: a nonempty groupoid; a commutative monoid, a group, a monoid with zero; a finite monoid with a single minimal right ideal; a category with an initial object, such as $\Delta^{\circ}$, the augmented simplex category, or finite sets with the empty set; the free category on a quiver in which every vertex has at most one incoming arrow (\lean{rightOreCondition_paths_iff}). None is a local test category among: posets with a minimal element, since the slice over it is the point; finite free categories on acyclic quivers, for the same reason at a source.
\end{corollary}

For finite monoids, the Ore condition is the existence of a single minimal right ideal: two minimal right ideals are disjoint, and every right ideal $mM$ contains one. For the second list we use the last assertion of Theorem~\ref{th:ore} with $X$ representable.

\status{} Grothendieck states the test-category case for an initial object in \emph{Pursuing Stacks}: \q{it is easy to see that a test category cannot possibly have an initial object} \cite[\S88, typescript p.~294]{ps-arxiv}, without an argument; Maltsiniotis notes that a representable separating segment excludes an initial object \cite[Rem.~6.8]{maltsiniotis-2001}, and the families of test categories of \cite[Ex.~4.1.20]{cisinski-308} and \cite[Ex.~6.9]{maltsiniotis-2001} exclude the empty category. So the answer \q{non} of folder 108 to question~6 agrees with Grothendieck's own later remark. What we did not find in these sources, nor through TheoremSearch and alphaXiv, is the general form: the Ore condition, that is, the De Morgan property of the topos, rules out weak and local test categories. The bridge reads it in one line: a logical property of $\widehat A$ (the De Morgan law) forbids a homotopical one (a contractible object of truth values). The converse fails: $\Delta$ does not satisfy the Ore condition and is a test category, while the non-Ore finite monoids below are not.

\subsection{What the arches generate}

Composed with the arches already proved, Theorem~\ref{th:ore} gives three families of statements.

\emph{Ore decks with no point slice.} The elementary argument of Corollary~\ref{cor:ore} needs an object with no incoming arrow but its identity; the Ore condition does not. By the preorder form of the De Morgan arch (\lean{rightOreCondition_preorder_iff}), a nonempty preorder in which every $\downarrow x$ is down-directed is neither a weak nor a local test category: so are $\mathbb Z$, $\mathbb Q$, $\mathbb R$ and every total order, although none of their slices is the point. The same holds for the monoid $\mathbb N$, whose presheaves are the sets with an endomorphism (not a local test category, while $B\mathbb N\simeq S^1$ already rules out a test category), and for the free category on a quiver in which every vertex has one incoming arrow and none is a source, such as $\cdots\to-1\to0\to1\to\cdots$. Among the 2510 monoids of order at most~6 of the atlas, 2015 satisfy the Ore condition and fall under the theorem with no homology computed.

\emph{Test categories are not De Morgan.} Read backwards, Theorem~\ref{th:ore} says that the presheaf topos of a weak or local test category is never De Morgan. For simplicial sets this was proved directly (\lean{not_isDeMorganTopos_simplicialSets}); it now follows for cubical sets \cite[Cor.~8.4.13]{cisinski-308}, for $\Theta_n$-sets and $\Theta$-sets \cite[Th.~6.5]{hubert-dk}, and for every test category recalled in \cite[introduction]{hubert-dk}, the dendroidal category included.

\begin{proposition}\label{prop:slices}
Let $W$ be a non-coarse fundamental localizer and $A$ a local $W$-test category. For every nonempty presheaf $X$ on $A$, the slice topos $\widehat A/X\simeq\widehat{A/X}$ is not De Morgan. In particular, for every object $a$ some object with an arrow to $a$ is the target of a cospan that does not complete to a commutative square.
\end{proposition}

\begin{proof}
$A/X$ is a local test category \cite[Prop.~4.1.14 (v)]{cisinski-308}, nonempty since $X$ is; by Theorem~\ref{th:ore} it does not satisfy the Ore condition, and by the De Morgan arch $\widehat{A/X}$ is not De Morgan. For $X = a$, a cospan of $A/a$ is a cospan of $A$ whose apex maps to $a$, and a commutative square in $A/a$ is one in $A$.
\end{proof}

So a test topos is not De Morgan anywhere, not only globally. We did not search the literature for Proposition~\ref{prop:slices}; it is a two-line consequence of the theorem and of Cisinski's Proposition~4.1.14.

\subsection{Finite decks: computations and a conjecture}

On the decks the generator enumerates, the script \nolinkurl{computations/187-test-categories/finite_test.py} computes the homology in degrees $0$ to $2$, over $\mathbb Q$ and $\mathbb F_2$, of the nerves of $A$, $A/L_A$ and every $A/(a\times L_A)$. For every monoid of order at most~6 (2510 up to isomorphism, from the atlas of the companion article), the homology of $A/L_A$ in degrees at most~2 differs from that of a point, and so does that of $A/(a\times L_A)$: no monoid of order $\le 6$ is a weak test category or a local test category, for the usual weak equivalences. The non-Ore ones, 495 of them, are not covered by Theorem~\ref{th:ore}; for them $H_1(A/(a\times L_A);\mathbb Q)$ has rank between $4$ and $743$. For posets, \nolinkurl{posets_euler.py} computes the Euler characteristic of the finite poset $A/L_A$ by the formula $\chi(A/L_A) = \sum_x |D(\downarrow x)|\,c(x)$, with $D(\downarrow x)$ the down-sets of $\downarrow x$ and $c(x) = 1-\chi(A_{<x})$: on every poset with at most six points and $\chi(A)=1$, $\chi(A/L_A)\ne1$, so no poset with at most six points is a weak test category (a poset with $\chi(A)\ne1$ is not even aspherical).

\begin{remark}[Conjecture]
No nonempty finite category is a weak test category, nor a local test category. The local case holds for finite posets and finite acyclic free categories (Corollary~\ref{cor:ore}); for monoids we only have the computations above.
\end{remark}

\subsection{Infinite products}

\begin{proposition}\label{prop:products}
For every non-coarse fundamental localizer $W$, an infinite product of $W$-aspherical categories need not be $W$-aspherical, and the class $W$ is not stable under infinite products.
\end{proposition}

\begin{proof}
Let $Z_n$ be the zigzag $0\to1\leftarrow2\to\cdots$ with $n+1$ objects, a contractible poset. A zigzag of length $k$ in $\prod_n Z_n$ projects onto zigzags of length $k$ in every factor, and the ends of $Z_n$ are at distance $n$; so $(0)_n$ and $(n)_n$ lie in different components of $\prod_n Z_n$, which is not $W$-aspherical (Theorem~\ref{th:ore}, proof). Hence $\prod_n(Z_n\to e)$ is not in $W$.
\end{proof}

This answers the question of folder 118, p.~36, which Grothendieck there judges less important than infinite sums; the argument is that of p.~28, where the same zigzags show that $(\mathrm{Cat})\to\overline{\mathrm{Cat}}$, the passage to homotopy classes of functors, does not commute with infinite sums. Since the nerve commutes with limits, it also gives an infinite product of aspherical simplicial sets that is not aspherical.

\leanentry{} Five files of \texttt{lean/bridges}. \lean{TestCategories.lean} proves the connectedness core with Mathlib's Lawvere object \lean{Functor.sieves}: Lemma~\ref{lem:ore-sieves} (\lean{rightOreCondition_iff_pullback_ne_bot}), the splitting of every category of elements over $L_A$ (\lean{not_isPreconnected_elements}) and its consequences (\lean{not_isConnected_elements_sieves}, \lean{not_isConnected_elements_yonedaProdSieves}, \lean{not_isConnected_elements_sieves_of_isDeMorganTopos}, \lean{not_isDeMorganTopos_of_isConnected_elements_sieves}). Mathlib has neither the weak equivalences of $\mathrm{Cat}$ nor Grothendieck's test notions, and neither has the Stacks library of AxiomMath, which we checked; so \lean{FundamentalLocalizers.lean} states them axiomatically, as Grothendieck did: a fundamental localizer (\lean{FundamentalLocalizer}: weak saturation, categories with a terminal object aspherical, relative Theorem~A), $i_A^*$ (\lean{iStar}), weak and local test categories (\lean{IsWeakTest}, \lean{IsLocalTest}, after Déf.~4.1.8 and~4.1.12 of \cite{cisinski-308}). Under the hypothesis that aspherical categories are connected (\lean{AsphericalConnected}, which is what $W\subset W_0$ gives), it proves Theorem~\ref{th:ore} in full (\lean{not_isWeakTest}, \lean{not_isLocalTest}, \lean{not_isLocalTest_of_over}), its reading through the De Morgan arch and backwards (\lean{not_isWeakTest_of_isDeMorganTopos}, \lean{not_isDeMorganTopos_of_isWeakTest}), Proposition~\ref{prop:slices} for every presheaf $X$ with a section (\lean{not_isDeMorganTopos_slice_of_isLocalTest}, first for $X$ representable, \lean{not_isDeMorganTopos_over_of_isLocalTest}), without Cisinski's Proposition~4.1.14: if $A/X$ satisfied the Ore condition, so would its slice over an object $p=(a,x)$, then, along a section of the projection $(A/X)/p\to A/a$ up to isomorphism (\lean{overCostructuredArrowSection}, \lean{rightOreCondition_of_iso}), the weak test category $A/a$, and for the slice topos $\widehat A/X$ itself through Mathlib's equivalence $\widehat A/X\simeq\widehat{A/X}$ (\lean{overEquivPresheafCostructuredArrow}; \lean{isDeMorganTopos_over_iff}, \lean{not_isDeMorganTopos_overPresheaf_of_isLocalTest}), the corollaries for $\Delta^{\circ}$, groupoids, down-directed preorders and $\mathbb Z$, the monoid $\mathbb N$, finite monoids with a single minimal right ideal, free categories with at most one arrow into each vertex, preorders with a minimal element, and finite free categories without cycles, through the existence of a source (\lean{not_isTest_simplexCategory_op}, \lean{not_isTest_groupoid}, \lean{not_isTest_preorder}, \lean{not_isTest_int}, \lean{not_isTest_multiplicative_nat}, \lean{not_isTest_finite_monoid}, \lean{not_isTest_paths}, \lean{not_isLocalTest_of_minimal}, \lean{exists_source_of_acyclic}, \lean{not_isLocalTest_paths_of_acyclic}), and Proposition~\ref{prop:products} in full: each zigzag $Z_n$ is aspherical for every fundamental localizer, from the axioms alone (\lean{isAspherical_zigZag}: induction on $n$, through the relative Theorem~A applied to the inclusion $Z_n\to Z_{n+1}$ or to the retraction $Z_{n+1}\to Z_n$, whose slices have a terminal object), while their product is not connected, hence not aspherical (\lean{not_isPreconnected_zigZagProd}, \lean{isAspherical_zigZag_and_not_isAspherical_zigZagProd}). Not formalised: that the usual weak equivalences, through the nerve, form a fundamental localizer contained in $W_0$, and Cisinski's Proposition~9.3.2 (both need the weak equivalences of simplicial sets); and the test status of $\Theta$ and $\Omega$. \lean{Theta2.lean} defines Joyal's $\Theta_2=\Delta\wr\Delta$, which no library has (\lean{Theta2}: a morphism is a map of $\Delta$ on the vertices with, for each cell of the target over a cell of the source, a map of $\Delta$ between their two-cells), and proves that it is aspherical and does not satisfy the Ore condition (\lean{Theta2.isAspherical}, \lean{Theta2.not_rightOreCondition}), so the obstruction of Theorem~\ref{th:ore} does not apply to it. Its test status is not formalised. The cone method fails: a final one-cell would need a vertex in each $[b_j]$ beyond $\varphi(m)$, and no choice is natural. The cylinder method fails too: thickening the two-cells is a functor with two natural ends, but a map that collapses a cell forgets the components over it, so the contraction is not natural. The known proofs go through Grothendieck's criterion \cite[Th.~4.1.19]{cisinski-308} or through the stability of test categories under $\Delta\wr-$. \lean{SimplexTest.lean} proves, from the axioms alone and with no connectedness hypothesis, that $\Delta$ is a weak and a local test category for every fundamental localizer (\lean{isWeakTest_simplexCategory}, \lean{isLocalTest_simplexCategory}, \lean{isTest_simplexCategory}). Two lemmas carry the homotopy: a natural transformation between two functors makes one a $W$-equivalence iff the other is (\lean{mem_iff_of_natTrans}, through the projection $\Delta_1\times D\to D$ and the relative Theorem~A), and a category with an endofunctor $t$ and natural transformations $1\to t\leftarrow\mathrm{const}_d$ is aspherical (\lean{isAspherical_of_contraction}). For $C$ with a final object $\star$, $A/i_A^*(C)$ is contracted by $(a,F)\mapsto(Ta,F^+)$, where $Ta$ adds a final vertex to $a$ and $F^+$ is $F$ on the simplices of $Ta$ that miss the new vertex and $\star$ on the others (\lean{ConeData}, \lean{ConeData.isWeakTest}); $\Delta$ and each slice $\Delta/b$ carry this cone structure (\lean{coneData}, \lean{coneDataOver}, the new vertex lying over the top of $b$). \lean{CubeTest.lean} does the same for the cube category $\square$ (faces and degeneracies, \lean{Cube}; \lean{isTest_cube}, \cite[Cor.~8.4.13]{cisinski-308} for the usual weak equivalences). $\square$ has no natural cone vertex but has a cylinder $[1]^n\mapsto[1]^{n+1}$ with two ends; the end $1$ gives a natural transformation from $(a,F)\mapsto(a,\star)$ to the contraction, so the projection $A/i_A^*(C)\to A$ is a $W$-equivalence, a retraction of $a\mapsto(a,\star)$, and $A/i_A^*(C)$ is aspherical since $A$ has a final object (\lean{CylData}, \lean{CylData.isWeakTest}); each slice $\square/b$ carries the cylinder over $b$ through the projection (\lean{cylDataOver}).
